\SetPractica{Ácidos carboxílicos}{Obtención e identificación de ácidos carboxílicos}

\section{Introducción teórica}



\section{Objetivos}

\begin{itemize}
    \item Obtener por un proceso de destilación ácido etanoico e identificarlo posteriormente a través de una prueba química.
\end{itemize}

\section{Materiales, equipos y reactivos}

\begin{table}[H]
\centering{\small\setstretch{1.25}
\begin{tabular}{|m{0.175\linewidth}|m{0.225\linewidth}|m{0.45\linewidth}|} \hline
    
    \thead{Equipos} & \thead{Materiales} & \thead{Reactivos} \\ \hline
    
    {Balanza \par}
    {Equipo de destilación}
    
    &
    
    {Pipeta graduada \par}
    {Soportes universales \par}
    {Mechero bunsen \par}
    {Soporte para mechero \par}
    {Malla de asbesto \par}
    {Probeta \par}
    {Tubos de ensayo \par}
    {Pipetas pasteur \par}
    {Espátulas}
    
    & 
    
    {Ácido sulfúrico (\ce{H2SO4}) concentrado \par}
    {Acetato de sodio (\ce{CH3COONa}) \par}
    {Solución acuosa de yodato de potasio (\ce{KIO3}) al 4\% \par}
    {Solución acuosa de yoduro de potasio (\ce{KI})   al 2\% \par}
    {Ácido acético (\ce{CH3COOH}) \par}
    {Solución acuosa de almidón al 0.1\% \par}
    {Vinagre blanco}
    
    \\ \hline
\end{tabular}}\end{table}

\section{Procedimientos}

\subsection{Preparación de solución de dicromato de potasio al 5\%}
\begin{enumerate}
    \item Pesar cuidadosamente 1.5 g de dicromato de potasio (usar guantes y tapaboca)
    \item Disolverlos en 30 ml de agua destilada.
\end{enumerate}

\subsection{Obtención de ácido etanóico}

\begin{enumerate}
    \item Coloque 7.0 g de acetato de sodio en un matraz de destilación de 250 mL
    \item Arme el aparato de destilación de acuerdo a las instrucciones de su profesor.
    \item Vierta al matraz 5 ml de ácido sulfúrico concentrado a través del embudo de seguridad.
    \item Caliente la mezcla hasta ebullición y reciba el destilado en una probeta de 10 ml.
    \item Al término de la destilación apague el mechero y desconecte las mangueras. 
\end{enumerate}

\subsection{Identificación de ácidos carboxílicos (método yodato-yoduro)}

\begin{enumerate}
    \item Añadir en un tubo pequeño:
    \begin{enumerate}
        \item 2 gotas del destilado (ácido etanóico)
        \item 2 gotas de cada una de las muestras proporcionadas por el profesor.
    \end{enumerate}
    
    \item Añadir al tubo que contiene las muestras y el destilado:
    \begin{enumerate}
        \item 2 gotas de solución acuosa de KI al 2\%
        \item 2 gotas de KIO3 al 4\%.
    \end{enumerate}

    \item Añadir a otro tubo 2 gotas de ácido acético e iguales cantidades de yodato y yoduro de potasio.
    \begin{enumerate}
        \item Dicho tubo se utilizará como control.
    \end{enumerate}

    \item Tapar los tubos e introducirlos en un baño de agua a 100°C durante 1 minuto.
    \item Enfriar los tubos y añadir 2 gotas de solución reciente de almidón al 0.1\%
\end{enumerate}

\textbf{Nota:} En el vial control y en las muestras que sean ácidos orgánicos la reacción producirá un color azul.